\subsection{Greatest Common Factor}
\begin{definition}[Greatest Common Factor (integers)]
A \term{common factor} of two (or more) integers is a number that is a factor of both (or all). The \term{greatest common factor} (or \term{GCF}) of two (or more) integers is their largest (positive) common factor.
\end{definition}
\begin{example}
List all (positive) factors of the integers $8$ and $20$. Then, give their GCF.
\begin{itemize}
\item\makeblankfill{Factors of 8}
\item\makeblankfill{Factors of 20}
\end{itemize}
The largest number on both lists is \makeblanksmall, so this is the greatest common factor of 8 and 20.
\end{example}
\noindent We can also use a calculator or computer to find the GCF of two numbers.
\begin{definition}[Greatest Common Factor (monomials)]
A \term{common factor} of two (or more) monomials is a monomial that is a factor of both (or all). The \term{greatest common factor} (or \term{GCF}) of two (or more) monomial is their common factor that has:
\begin{itemize*}
\item the largest possible coefficient
\item the highest possible degree
\end{itemize*}
\end{definition}
To find the GCF of monomials, we find:
\begin{itemize*}
\item the GCF of the coefficients
\item the lowest power of each variable that appears in the monomials
\end{itemize*}
If we take the prime factorization of each coefficient, we can treat prime numbers like variables. Or, we can use a calculator to find the GCF of the coefficients.
\begin{example}
Find the GCF of $3x^4y^3$, $6x^3y^5$, and $9x^2y^7$.
\begin{itemize}
\item\makeblankfill{For the coefficient: }
\item\makeblankfill{For $x$: }
\item\makeblankfill{For $y$: }
\end{itemize}
The GCF of $3x^4y^3$, $6x^3y^5$, and $9x^2y^7$ is \makeblanksmall.
\end{example}